% shared/figures/tikz/_style.tex — ARI house design language for system figures.
% ---------------------------------------------------------------------------
% \input this file at the TOP of a figure source, BEFORE \begin{tikzpicture}.
% It only DEFINES colors / \tikzset styles / pics — it typesets nothing — so
% it is safe to \input from several figures in the same document and safe to
% load in any engine (pdflatex / lualatex / xelatex): xcolor + TikZ core only,
% no fontspec, no external icon packages (glyphs below are stroke-drawn pics).
%
% Relationship to the legacy styles: the baseline block/arrow/group styles in
% style.tikzstyles is loaded by shared/preamble.tex / the standalone wrapper;
% this file ADDS the refined house language (flat pastel phase fills, thin-ink
% rounded boxes, dashed loop enclosures, elbow handoffs) and never redefines a
% baseline style name. New-style figures should use ONLY the ari* styles.
%
% inline-width: vendor-coloured override — gate-T2 marker: this is a shared
% STYLE file. minimum width / minimum height / text width are centralised
% here (exactly as in style.tikzstyles) and must NOT be re-specified inline
% in figure sources.

% --------------------------------------------------------------- palette
% Flat pastels on white. FILL ENCODES PHASE (which lifecycle phase owns the
% box); outlines, arrows, headings and body text are a single soft near-black
% ink; gray is reserved for loop enclosures, notes and secondary flows;
% `accent` is the only saturated color and is rationed to at most ONE
% semantic role per figure (e.g. the blocking verdict edge).
\definecolor{phaseidea}{HTML}{DBE5F6}%   lavender-blue — ideation / planning
\definecolor{phaseexp}{HTML}{FBE3C9}%    peach         — experimentation / execution
\definecolor{phasepaper}{HTML}{E9E9E9}%  light gray    — writing / publication
\definecolor{phasegate}{HTML}{D9EDDC}%   pale green    — gates / evaluation / verdicts
\definecolor{accent}{HTML}{D55E00}%      vermilion     — the ONE emphasis color (Okabe-Ito)
\definecolor{ariink}{HTML}{3C3C3C}%      soft near-black — outlines, arrows, headings
\definecolor{arimuted}{HTML}{8A8A8A}%    mid gray      — loop enclosures, notes, dashes

% --------------------------------------------------------------- styles
\tikzset{
  % Canvas rhythm — put on the picture itself: \begin{tikzpicture}[aricanvas].
  % Tight gaps INSIDE a phase column (0.5cm vertical), generous whitespace
  % BETWEEN phase columns (1.6cm horizontal): the phase-grouped-column layout.
  aricanvas/.style={node distance=0.5cm and 1.6cm,
                    line cap=round, line join=round},
  % ----- boxes: rounded rectangle, thin ink outline, flat pastel fill ------
  % Default fill is white; phase membership is declared with the fill, either
  % via the shorthands below or explicitly: \node[aribox, fill=phaseexp] ...
  aribox/.style={rectangle, rounded corners=3pt, draw=ariink,
                 line width=0.6pt, fill=white,
                 minimum width=2.6cm, minimum height=0.8cm, text width=2.4cm,
                 align=center, inner sep=3pt, font=\sffamily\footnotesize,
                 text=ariink},
  ariboxwide/.style={aribox, minimum width=3.6cm, text width=3.4cm},
  aribox idea/.style={aribox, fill=phaseidea},
  aribox exp/.style={aribox, fill=phaseexp},
  aribox paper/.style={aribox, fill=phasepaper},
  aribox gate/.style={aribox, fill=phasegate},
  % ----- phase heading: bold sans ink text above each column ---------------
  ariphase/.style={font=\bfseries\sffamily\small, text=ariink,
                   align=center, inner sep=2pt},
  % ----- loop enclosure: dashed gray rounded rectangle, fit-positioned -----
  %   \node[ariloop, fit=(a)(b)] (loop) {};
  ariloop/.style={rectangle, draw=arimuted, dashed, line width=0.7pt,
                  rounded corners=7pt, inner sep=8pt},
  % ----- arrows -------------------------------------------------------------
  % ariarrow        primary flow (short, solid, ink)
  % ariarrow muted  secondary / data / observation flow (dashed gray)
  % ariarrow accent the rationed emphasis edge (solid vermilion)
  % arifeedback     long outer feedback path (dotted gray, rounded routing)
  ariarrow/.style={-{Stealth[length=5pt,width=4pt]}, draw=ariink,
                   line width=0.7pt, shorten <=2.5pt, shorten >=2.5pt},
  ariarrow muted/.style={ariarrow, draw=arimuted, dashed},
  ariarrow accent/.style={ariarrow, draw=accent},
  arifeedback/.style={ariarrow, draw=arimuted, densely dotted,
                      rounded corners=7pt},
  % ----- elbow handoff between phase columns --------------------------------
  % Used with `to`: \draw[arielbow] (lastA) to (firstB);
  % arielbow    goes horizontal first, then vertical  (-|)
  % arielbow vh goes vertical first, then horizontal  (|-)
  arielbow/.style={ariarrow, rounded corners=5pt,
                   to path={-| (\tikztotarget) \tikztonodes}},
  arielbow vh/.style={ariarrow, rounded corners=5pt,
                      to path={|- (\tikztotarget) \tikztonodes}},
  % ----- small annotation / sparse edge label --------------------------------
  arinote/.style={font=\sffamily\scriptsize, text=arimuted, align=center,
                  fill=white, inner sep=1.5pt, rounded corners=1pt},
  % ----- numbered step badge (replaces edge labels for flow ordering) -------
  aribadge/.style={circle, fill=ariink, text=white, inner sep=1pt,
                   minimum size=11pt, font=\bfseries\sffamily\scriptsize},
  % ----- stroke-drawn heading glyphs (no icon packages) ----------------------
  % ~3mm tall, designed to sit left of an ariphase heading:
  %   \node[ariphase] (h) {\figlabel{Fxx.head.idea}};
  %   \pic at ([xshift=-2.5mm]h.west) {ariglyph bulb};
  ariglyph bulb/.pic={% ideation (lightbulb)
    \draw[ariink, line width=0.6pt] (0,0.45mm) circle [radius=0.95mm];
    \draw[ariink, line width=0.6pt] (-0.45mm,-0.85mm) -- (0.45mm,-0.85mm);
    \draw[ariink, line width=0.6pt] (-0.32mm,-1.3mm) -- (0.32mm,-1.3mm);
  },
  ariglyph flask/.pic={% experimentation (Erlenmeyer flask)
    \draw[ariink, line width=0.6pt]
      (-0.35mm,1.55mm) -- (-0.35mm,0.25mm) -- (-1.15mm,-1.45mm)
        -- (1.15mm,-1.45mm) -- (0.35mm,0.25mm) -- (0.35mm,1.55mm);
    \draw[ariink, line width=0.6pt] (-0.65mm,1.55mm) -- (0.65mm,1.55mm);
  },
  ariglyph doc/.pic={% writing / publication (page with dog-ear)
    \draw[ariink, line width=0.6pt]
      (-1.05mm,1.5mm) -- (-1.05mm,-1.5mm) -- (1.05mm,-1.5mm)
        -- (1.05mm,0.65mm) -- (0.2mm,1.5mm) -- cycle;
    \draw[ariink, line width=0.6pt] (-0.55mm,0.1mm) -- (0.55mm,0.1mm);
    \draw[ariink, line width=0.6pt] (-0.55mm,-0.6mm) -- (0.55mm,-0.6mm);
  },
  ariglyph gate/.pic={% gate / evaluation (check in circle)
    \draw[ariink, line width=0.6pt] (0,0) circle [radius=1.4mm];
    \draw[ariink, line width=0.6pt]
      (-0.65mm,0.05mm) -- (-0.15mm,-0.55mm) -- (0.75mm,0.6mm);
  },
  ariglyph loop/.pic={% iteration marker, sits inside an ariloop enclosure
    \draw[ariink, line width=0.6pt, -{Stealth[length=1.5pt,width=1.4pt]}]
      (115:1.2mm) arc[start angle=115, end angle=-50, radius=1.2mm];
    \draw[ariink, line width=0.6pt, -{Stealth[length=1.5pt,width=1.4pt]}]
      (-65:1.2mm) arc[start angle=-65, end angle=-230, radius=1.2mm];
  },
}%
